\documentclass[11pt]{article}
\usepackage[margin=24mm]{geometry}
\usepackage{amsmath,amssymb,amsthm,booktabs,hyperref}
\newtheorem{theorem}{Theorem}
\title{Scalar bounds that hold at cutoff five}
\author{Research note}\date{7 October 2026}
\begin{document}\maketitle
\noindent\textbf{Status.} The statements below are proved at the requested
cutoff $5$. They are necessary conditions, not a proof of
$M(H)<5\Rightarrow d_F(H)<1/30$. In particular, the small-$J$ alternative
still needs to be converted into Fourier proximity or the final scalar
alternatives by using the rest of the reconstruction.

For a corner of a raw order-six complex Hadamard, use
$p=X+Y+Z$ and $J=2\operatorname{Re}(sr\bar t)$ in the usual complementary
Gram notation. Complementing its three columns preserves $p$ and replaces
$J$ by $-J$. Call the member with the larger $R$ the \emph{oriented member}
of that pair. Ties may be oriented arbitrarily.

\begin{theorem}\label{main}
Suppose $M(H)<5$. Then every corner satisfies
\[
 \boxed{|J|<\frac{21}{8}.}
\]
For an oriented member with $0<p<5$, one has
\[
 \boxed{J\ge0,\qquad 12Jp(9-p)(5-p)<25.}
\]
Consequently the whole region $1\le p\le4$ satisfies
\[
 \boxed{0\le J<\frac5{48}.}
\]
The only finite certificate in this proof consists of 24 rational
Bernstein comparisons, displayed below as integer upper bounds.
\end{theorem}

\paragraph{The complementary-corner identity.}
Write $t=9-p$ (now a scalar, not a Gram entry), $z=J^2$, and
$z_0=t^3(4-t)/4$. For an oriented pair put
\[
 U=\frac{R_++R_-}{2},\qquad V=\frac{R_+-R_-}{2}.
\]
The reconstruction gives
\begin{equation}\label{id}
 4z(z-z_0)(z+36p)
 =JV(2z+30t^2)-Uz(t+6)^2-\frac{V^2}{4}.
\end{equation}
Also $0\le V\le U$, $U+V<5$, and $0\le |J|\le t\le9$.
The first two facts follow from nonnegativity and orientation of the two
discriminants. The last follows from the complementary permanent identity
$p+|J|\le9$. Identity \eqref{id} is a polynomial expansion of those
discriminants, verified in \texttt{complementary\_scalar\_rule.py}.
It remains valid when both discriminants are zero.

If $J<0$, the right side is nonpositive, hence
$0<J^2\le z_0$. This forces $0<t<4$, or $5<p<9$, and gives
\[
 J^2\le\frac{27}{4},\qquad
 27-t^3(4-t)=(t-3)^2(t^2+2t+3)\ge0.\tag{1}
\]
In particular, the oriented member cannot have $J<0$ when $p<5$.

If $J>0$ and $p\le7$, define
\[
 D=2J^2+30t^2-J(t+6)^2,\qquad
 Q(p,J)=5D-8J(J^2+36p)(J^2-z_0).
\]
Here $2\le t\le9$, and
\[
 \frac D J=\frac{2(J-t)(J-15t)}J+(t-2)(18-t)\ge0.
\]
Dividing \eqref{id} by $J$, using $U\ge V$, and then $V<5/2$ gives
\[
 4J(J^2-z_0)(J^2+36p)\le VD\le\frac52D,
 \qquad\text{so }Q(p,J)\ge0.\tag{2}
\]

\paragraph{The region $0<p<5$.}
We already know $J\ge0$. If $J=0$, the assertion is immediate.
Suppose $J>0$ and put $L=Jp\,t(t-4)$. Here $4<t<9$ and
$t(t-4)<45$. Expanding $Q$ in a useful grouping gives
\begin{align*}
 Q={}&t^2(150-72L)-5J(t+6)^2\\
 &+2J^2\bigl(5-4J^3-144pJ\bigr)-2J^3t^3(t-4).
\end{align*}
If $L\ge25/12$, the first term is nonpositive, and
\[
 144pJ=\frac{144L}{t(t-4)}>\frac{300}{45}>5.
\]
All other displayed terms are strictly negative. This contradicts (2).
Thus $12L<25$. On $1\le p\le4$,
\[
 p(9-p)(5-p)-20=(4-p)(-p^2+10p-5)\ge0,
\]
because $-p^2+10p-5=4+(p-1)(9-p)\ge4$.
Therefore $240J<25$, proving $J<5/48$.

\paragraph{The bound on $|J|$.}
It suffices to work with the oriented member. Negative $J$ is covered
by (1), since $27/4<(21/8)^2$. For positive $J$, suppose
$J\ge a=21/8$. The elementary estimate
\[
 J^2\le4XYZ\le4(p/3)^3
\]
gives $p>3$, hence $t<6$. Also $J\le t$ implies $p\le51/8<7$,
so (2) applies.

If $t\ge4$, then $z_0\le0$ and $D\le32t^2<1152$, while
\[
 8J(J^2+36p)(J^2-z_0)
 >8\left(\frac52\right)108\left(\frac{25}{4}\right)
 =13500>5760>5D.
\]
This contradicts (2). It remains to consider $a\le t\le4$, or
$5\le p\le51/8$.
For fixed $p$ on this interval, $D$ decreases as $J$ increases through
$[a,t]$, since $4J-(t+6)^2<0$. Each positive factor of
$8J(J^2+36p)(J^2-z_0)$ increases, since
$J^2\ge a^2>27/4\ge z_0$.
Consequently $Q(p,J)\le Q(p,a)$.
The following exact certificate proves $Q(p,a)<0$ throughout the interval,
again contradicting (2). This proves the theorem.

\paragraph{The 24-number certificate.}
On each row's interval $[l,h]$, write
$Q(l+(h-l)u,21/8)$ in the degree-five Bernstein basis
$\binom5k u^k(1-u)^{5-k}$, $0\le k\le5$.
The table gives integer \emph{upper bounds} for its six exact rational
coefficients. Every entry is negative, and the basis is nonnegative and
sums to one.
\begin{center}\small
\begin{tabular}{c rrrrrr}
\toprule
$[l,h]$ & $0$ & $1$ & $2$ & $3$ & $4$ & $5$\\\midrule
$[5,171/32]$ & $-25887$ & $-21992$ & $-18510$ & $-15428$ & $-12733$ & $-10407$\\
$[171/32,91/16]$ & $-10407$ & $-8081$ & $-6124$ & $-4519$ & $-3246$ & $-2283$\\
$[91/16,193/32]$ & $-2283$ & $-1320$ & $-668$ & $-304$ & $-205$ & $-346$\\
$[193/32,51/8]$ & $-346$ & $-488$ & $-870$ & $-1469$ & $-2258$ & $-3214$\\
\bottomrule
\end{tabular}
\end{center}
The script \texttt{practical\_scalar\_regions.py} regenerates these values
using rational arithmetic and independently reconstructs the entire
polynomial from each set of six coefficients. Its JSON receipt includes
the exact fractions and source hashes. No sampling or floating-point
comparison is used in this certificate.
\end{document}
